\section{Discussion and Conclusion}
\label{sec:conclusion}

This thesis asks whether the Variance Risk Premium creates more economic value as an
informational state variable for dynamic equity--bond allocation or through a model-based direct
variance-payoff approximation, and whether the answer is stable across the United States and
Europe.

The empirical answer is conditional. VRP-related variables contain information about market
stress, but adding them to HMM, Markov-switching, or machine-learning allocation models does
not generate robust benchmark dominance. The direct-payoff channel produces stronger evidence,
particularly in Europe. The main conclusion is therefore not that VRP is universally profitable,
but that its economic value depends critically on the mechanism through which variance
information is transformed into portfolio returns.

\subsection{Four implications for portfolio management}

\textbf{1. Do not treat VRP as a universal allocation alpha.}

In the United States, the strongest HMM, HMM RV + Log VRP, reaches a Sharpe ratio of
1.019, compared with 1.025 for the equal-weighted equity--bond benchmark. In Europe, the
strongest regime model reaches 0.281, below buy-and-hold equity at 0.313. Machine-learning
results lead to the same conclusion: VRP improves selected predictive diagnostics, particularly
for the European Random Forest, but this does not translate into statistically significant welfare
dominance.

\textbf{Do:} use VRP as a supplementary state variable when it improves economically relevant
risk information.

\textbf{Do not:} infer from statistical predictability that a VRP-enhanced allocation rule should
replace simple diversified benchmarks.

\medskip

\textbf{2. Use regime models primarily as risk-shaping overlays rather than return engines.}

The strongest practical contribution of the United States HMM is downside-risk management
rather than higher average risk-adjusted return. Its maximum drawdown is -16.67\%, compared
with -20.06\% for 60/40 and -19.10\% for equal weighting. Full monthly implementation,
however, generates average turnover of 25.35\%. Partial rebalancing with an adjustment speed of
0.25 reduces turnover to 9.10\% while preserving most of the risk characteristics.

\textbf{Do:} translate regime probabilities gradually into portfolio exposure and evaluate the
model through drawdown, turnover, and welfare as well as Sharpe ratio.

\textbf{Do not:} fully trade every monthly change in estimated state probability when the gross
performance advantage over static portfolios is small.

\medskip

\textbf{3. The strongest economic evidence is in variance carry, but only under a clearly defined
modeling boundary.}

The selected United States High-VRP direct strategy produces a Sharpe ratio of 0.882, but at
risk aversion $\gamma=5$ its mean--variance certainty-equivalent differences are -70.9 basis
points relative to 60/40 and -45.0 basis points relative to equal weighting. The corresponding
bootstrap intervals include zero. The United States evidence therefore supports a defensive
variance-carry interpretation, not robust welfare superiority.

Europe is materially different. The VRP-positive direct-payoff strategy produces an annualized
return of 13.08\%, volatility of 9.70\%, and a Sharpe ratio of 1.324. Its mean--variance
certainty-equivalent advantages at $\gamma=5$ are approximately 927 basis points relative to
60/40 and 899 basis points relative to equal weighting, with positive lower bootstrap bounds.
The conclusion also survives the principal robustness tests, including transaction costs up to
50 basis points, alternative risk-estimation windows, a maximum notional reduced to 5\%,
subperiod analysis, crisis exclusions, and alternative risk-aversion coefficients.

\textbf{Do:} interpret the European result as evidence that the variance-carry payoff mechanism
deserves further investigation with contract-level derivative data.

\textbf{Do not:} interpret the reported 13.08\% return as the historical return of an executable
variance-swap strategy.

\medskip

\textbf{4. Do not transfer a VRP strategy mechanically across markets.}

The United States and European results are not interchangeable. VRP improves different model
families in the two markets, allocation performance differs materially, and the direct-payoff
evidence is substantially stronger in Europe. This heterogeneity may reflect differences in sample
history, volatility-index construction, market structure, crisis exposure, and measurement.

The allocation and direct-payoff layers also use different aligned samples: 184 and 122
out-of-sample allocation observations in the United States and Europe, compared with 232 and
170 direct-strategy observations. The thesis therefore does not interpret the two channels as a
single pooled horse race.

\textbf{Do:} validate variance signals and payoff rules separately by market and on matched
implementation assumptions.

\textbf{Do not:} infer a permanent structural European advantage from the historical
cross-market difference observed here.

\subsection{Interpretation boundary and next research step}

The strongest empirical result also carries the most important limitation. The direct-payoff
layer is a model-based approximation. Lagged squared VIX or VSTOXX is used as an
implied-variance strike proxy and trailing realized variance as settlement. The resulting
normalized payoff is then mapped into portfolio capital through lagged volatility targeting and a
notional constraint. It is not an observed return on invested capital.

A contract-level variance-swap reconstruction would require a sufficiently deep historical
option surface or institutional variance-swap quotations, together with maturity matching,
forward estimation, strike integration, bid--ask information, contractual notional conventions,
collateral, margin, and intramonth mark-to-market. Such data were not obtained and validated
within the empirical scope of this thesis. The direct-payoff evidence is therefore deliberately
reported as a model-based approximation rather than as executable historical variance-swap
performance.

Two additional limitations remain relevant. First, the number of independent extreme
volatility episodes is small, particularly for a negatively skewed payoff. Second, multiple HMM,
RSM, machine-learning, VRP, and robustness specifications are examined. Rolling
out-of-sample estimation, simple benchmarks, common-sample comparisons, and bootstrap
inference reduce specification-search risk but cannot eliminate it.

The highest-priority extension is consequently not another allocation model. It is to reconstruct
contract-matched forward variance from historical option surfaces and test the direct-payoff
conclusion on a genuinely tradable derivative representation. A later untouched holdout sample
and stronger multiple-model inference would provide the next layer of validation.

\subsection{Final answer}

The economic value of the Variance Risk Premium cannot be separated from the payoff,
benchmark, and implementation mechanism used to monetize it. As an informational variable,
VRP changes estimated states and selected predictive diagnostics but provides only limited
incremental value for dynamic equity--bond allocation. As a direct variance-carry input, it
produces substantially stronger evidence in the tested European sample.

The practical conclusion is therefore selective rather than universal: use VRP information as
a risk-conditioning input rather than as a guaranteed allocation alpha; implement regime
changes gradually when turnover is costly; treat the European direct-payoff result as a strong
research signal rather than an immediately executable return; and require market-specific
validation before deploying any VRP strategy.

In the United States, neither channel establishes robust investor-welfare dominance over
traditional portfolios. In Europe, the model-based direct variance-payoff approximation provides
the strongest evidence of economic value, subject to the derivative-data and implementation
limitations stated above.
